% !PRINTED 343
\weberpage{343}
If, however, $(1, \theta)$ is a basis of $\Omega$ and if, in order to comprise
both cases of §\,90, (9),
\weq{2}{\theta = \frac{-\Delta+\sqrt{\Delta}}{2}}
is assumed, then in case a)
\[
\theta \text{ not congruent with a number } r,
\]
in case b) or c)
\[
\theta \equiv r \pmod{\mathfrak{p}},
\]
so that, if one sets $-x = 2r + \Delta$, it follows:

1. \gesperrt{The necessary and sufficient condition for
the occurrence of one of the cases b), c) is this, that
there is a rational integer which is at once
even or odd with $\Delta$, and for which}
\weq{3}{\frac{x+\sqrt{\Delta}}{2} \equiv 0 \pmod{\mathfrak{p}}.}
If $\mathfrak{p}'$ is the ideal conjugate to $\mathfrak{p}$, then
\weq{4}{\frac{-x+\sqrt{\Delta}}{2} \equiv 0 \pmod{\mathfrak{p}'},}
and from this it follows
\weq{5}{x^2 \equiv \Delta \pmod{4p}.}
If conversely the congruence (5) is satisfied by $x = r$, then
\[
\frac{r+\sqrt{\Delta}}{2}\cdot\frac{-r+\sqrt{\Delta}}{2} \equiv 0 \pmod{\mathfrak{p}},
\]
and consequently either $x = r$ or $x = -r$ satisfies condition (3).

The two ideals $\mathfrak{p}, \mathfrak{p}'$ will be identical with one another
when the two numbers (3) and (4), and consequently their
sum $\sqrt{\Delta}$ as well, are divisible by $\mathfrak{p}$, and hence $\Delta$ by $p$. Then
$x \equiv 0$ or $x \equiv p \pmod{2p}$ is the only root of (5).

We thus have the result:

2. \gesperrt{The prime factor $\mathfrak{p}$ of the natural prime $p$ is
of the second degree when the congruence (5) has no
solution, of the first degree when (5) has a solution,
and $p$ is the square of $\mathfrak{p}$ when (5) has only
one solution, that is, when $p$ divides $\Delta$.}

The latter is in agreement with the general
theorem Bd.\,II, §\,174.
